\section*{Capítulo \ref{ch:g9:sqrt} --- Raízes quadradas}

\begin{solution}{exo:g9:sqrt:1}
$\sqrt{64} = 8$; \quad $\sqrt{100} = 10$; \quad
$\sqrt{0.09} = 0.3$ (pois $0.3^2 = 0.09$); \quad
$\sqrt{\frac{1}{25}} = \frac15$; \quad
$\left(\sqrt{13}\right)^2 = 13$; \quad
$\sqrt{(-4)^2} = \sqrt{16} = 4 = \abs{-4}$.
\end{solution}

\begin{solution}{exo:g9:sqrt:2}
$\sqrt{50} = \sqrt{25 \times 2} = 5\sqrt2$; \quad
$\sqrt{45} = \sqrt{9 \times 5} = 3\sqrt5$; \quad
$\sqrt{98} = \sqrt{49 \times 2} = 7\sqrt2$; \quad
$\sqrt{300} = \sqrt{100 \times 3} = 10\sqrt3$.
\end{solution}

\begin{solution}{exo:g9:sqrt:3}
$\sqrt2 \times \sqrt{18} = \sqrt{36} = 6$.

$\dfrac{\sqrt{75}}{\sqrt3} = \sqrt{\dfrac{75}{3}} = \sqrt{25} = 5$.

$\sqrt5 \times \sqrt{20} = \sqrt{100} = 10$.

$\dfrac{\sqrt8}{\sqrt{50}} = \sqrt{\dfrac{8}{50}} =
\sqrt{\dfrac{4}{25}} = \dfrac25$.
\end{solution}

\begin{solution}{exo:g9:sqrt:4}
$x^2 = 36$: $x = 6$ ou $x = -6$.

$x^2 = 11$: $x = \sqrt{11}$ ou $x = -\sqrt{11}$.

$x^2 + 4 = 0$ significa $x^2 = -4 < 0$: nenhuma solução.

$2x^2 = 50$: $x^2 = 25$, então $x = 5$ ou $x = -5$.
\end{solution}

\begin{solution}{exo:g9:sqrt:5}
$\sqrt{20} + \sqrt{45} = 2\sqrt5 + 3\sqrt5 = 5\sqrt5$.

$3\sqrt8 - \sqrt{18} + \sqrt2 = 3 \times 2\sqrt2 - 3\sqrt2 + \sqrt2
= (6 - 3 + 1)\sqrt2 = 4\sqrt2$.
\end{solution}

\begin{solution}{exo:g9:sqrt:6}
$\left(\sqrt3 + 1\right)^2 = 3 + 2\sqrt3 + 1 = 4 + 2\sqrt3$.

$\left(2\sqrt5 - 3\right)^2 = 4 \times 5 - 2 \times 3 \times 2\sqrt5
+ 9 = 29 - 12\sqrt5$.

$\left(\sqrt6 + \sqrt2\right)\left(\sqrt6 - \sqrt2\right) = 6 - 2 = 4$.
\end{solution}

\begin{solution}{exo:g9:sqrt:7}
Lado: $\sqrt{6400} = 80$ m. Diagonal (Pitágoras):
$\sqrt{80^2 + 80^2} = 80\sqrt2 \approx 113$ m.
\end{solution}

\begin{solution}{exo:g9:sqrt:8}
$\dfrac{1}{\sqrt2} = \dfrac{1 \times \sqrt2}{\sqrt2 \times \sqrt2}
= \dfrac{\sqrt2}{2}$.

$\dfrac{6}{\sqrt3} = \dfrac{6\sqrt3}{3} = 2\sqrt3$.
\quad
$\dfrac{10}{\sqrt5} = \dfrac{10\sqrt5}{5} = 2\sqrt5$.
\end{solution}

\begin{solution}{exo:g9:sqrt:9}
\emph{1. Verdadeiro}: é a regra do \omterm{def:g2:mult:def}{produto}
(\cref{thm:g9:sqrt:rules}).

\emph{2. Falso}: $\sqrt{9 + 16} = 5$, mas $\sqrt9 + \sqrt{16} = 7$.

\emph{3. Falso} para $x$ negativo: $\sqrt{(-4)^2} = 4 \neq -4$. O
enunciado correto é $\sqrt{x^2} = \abs{x}$.

\emph{4. Verdadeiro}: $\left(3\sqrt2\right)^2 = 9 \times 2 = 18$.
\end{solution}

\begin{solution}{exo:g9:sqrt:10}
$\left(2 + \sqrt3\right)^2 = 4 + 4\sqrt3 + 3 = 7 + 4\sqrt3$. Então
$x = \sqrt{7 + 4\sqrt3} = \sqrt{\left(2+\sqrt3\right)^2} = 2 + \sqrt3$
(um número não negativo, então a raiz apenas retira o quadrado).

Do mesmo modo, $\left(\sqrt5 - 2\right)^2 = 5 - 4\sqrt5 + 4 =
9 - 4\sqrt5$, e $\sqrt5 - 2 > 0$, então
$\sqrt{9 - 4\sqrt5} = \sqrt5 - 2$ (e não $2 - \sqrt5$, que é
negativo!).
\end{solution}

\begin{solution}{pb:g9:sqrt:1}
\textbf{1.} $1.4^2 = 1.96 < 2 < 2.25 = 1.5^2$;
$1.41^2 = 1.9881 < 2 < 2.0164 = 1.42^2$. Três casas decimais:
$1.414^2 = 1.999396 < 2 < 2.002225 = 1.415^2$, então
$1.414 < \sqrt2 < 1.415$.

\textbf{2.} Elevar $\sqrt2 = \frac ab$ ao quadrado dá
$2 = \frac{a^2}{b^2}$, logo $a^2 = 2b^2$: o número $a^2$ é o dobro de
um número inteiro --- \emph{par}.

\textbf{3.} Se $a$ fosse \omterm{def:g3:numbers:evenodd}{ímpar}, $a^2$ seria \omterm{def:g3:numbers:evenodd}{ímpar}
(\cref{pb:g8:pythagoras:1}); como $a^2$ é par, $a$ é par: $a = 2k$.
Substituindo: $4k^2 = 2b^2$, então $b^2 = 2k^2$ é par e, pelo mesmo
fato de paridade, $b$ é par.

\textbf{4.} $a$ e $b$ pares significa os dois divisíveis por $2$ ---
mas $\frac ab$ era irredutível, sem \omterm{def:g9:arith:divisor}{divisor} comum nenhum.
Contradição: a \omterm{def:g9:fractions:fraction}{fração} suposta não pode existir. \emph{Teorema:
$\sqrt2$ é irracional.}

\textbf{5.} A contradição mora inteirinha na palavra ``irredutível'':
sem ela, ``$a$ e $b$ os dois pares'' não contradiz coisa alguma. Supor
a \omterm{def:g9:fractions:fraction}{fração} irredutível é legítimo porque toda \omterm{def:g9:fractions:fraction}{fração} \emph{tem} uma forma
irredutível (\cref{met:g9:arith:simplify}: divida pelo MDC) --- se
$\sqrt2$ fosse uma \omterm{def:g9:fractions:fraction}{fração} qualquer, seria também uma \omterm{def:g9:fractions:fraction}{fração}
irredutível, e essa é impossível.

\textbf{6.} Nas decomposições em fatores primos, os quadrados carregam
\omterm{def:g8:powers:def}{expoentes} pares (\cref{pb:g9:arith:1}). Em $a^2 = 3b^2$, o \omterm{def:g8:powers:def}{expoente} do
$3$ é par à esquerda, mas \omterm{def:g3:numbers:evenodd}{ímpar} à direita (par, vindo de $b^2$, mais
um). Nenhum número tem duas decomposições
(\cref{thm:g9:arith:factorization}): contradição, e $\sqrt3$ é
irracional.

\textbf{7.} Em $a^2 = n b^2$, o argumento encontra um primo de
paridade de \omterm{def:g8:powers:def}{expoente} contraditória exatamente quando algum primo
aparece em $n$ com \omterm{def:g8:powers:def}{expoente} \emph{\omterm{def:g3:numbers:evenodd}{ímpar}}. Se todo \omterm{def:g8:powers:def}{expoente} de $n$ for
par, $n$ é um quadrado perfeito ($n = m^2$, e $\sqrt n = m$ é
inteiro). Resultado completo: $\sqrt n$ é irracional para todo inteiro
$n$ que não é quadrado perfeito ---
$\sqrt2, \sqrt3, \sqrt5, \sqrt6, \sqrt7, \sqrt8, \sqrt{10}, \dots$

\textbf{8.} De $r x = q$ com $r \neq 0$: $x = \frac qr$, um \omterm{def:g3:division:remainder}{quociente}
de racionais, logo racional. Então, se $x$ é irracional, $r x$ não pode
ser racional: $2\sqrt2$ e $\frac{\sqrt2}{2}$ são irracionais
(multiplicadores $2$ e $\frac12$).

\textbf{9.} Se $1 + \sqrt2 = q$ fosse racional, $\sqrt2 = q - 1$ seria
racional: contradição. Se $s = \sqrt2 + \sqrt3$ fosse racional, então
$s^2 = 2 + 2\sqrt6 + 3 = 5 + 2\sqrt6$ também seria, dando
$\sqrt6 = \frac{s^2 - 5}{2}$ racional --- mas $\sqrt6$ é irracional
(questão 7). Logo, $\sqrt2 + \sqrt3$ é irracional.

\textbf{10.} \omterm{def:g1:addition:def}{Somas}: $\sqrt2$ e $-\sqrt2$ (ou $\sqrt2$ e
$1 - \sqrt2$) são irracionais cada um, com \omterm{def:g1:addition:def}{somas} racionais $0$ e $1$.
\omterm{def:g2:mult:def}{Produtos}: $\sqrt2 \times \sqrt2 = 2$. A irracionalidade pode evaporar
nas operações --- daí as demonstrações caso a caso.

\textbf{11.} A partir de $x = 1$: média de $1$ e $\frac21 = 2$:
$x_1 = \frac32$. Depois, média de $\frac32$ e $\frac{2}{3/2} =
\frac43$: $x_2 = \frac12\left(\frac32 + \frac43\right)
= \frac12 \cdot \frac{17}{6} = \frac{17}{12} \approx 1.4167$.

\textbf{12.} $x_3 = \frac12\left(\frac{17}{12} + \frac{24}{17}\right)
= \frac12 \cdot \frac{289 + 288}{204} = \frac{577}{408} =
1.4142156\ldots$ Contra $\sqrt2 = 1.4142135\ldots$: cinco casas
decimais corretas depois de três passos --- a receita mais ou menos
dobra as casas corretas a cada rodada.

\textbf{13.} Se $x^2 > 2$ (chute grande demais), então $x > \sqrt2$ e
$\frac2x < \frac{2}{\sqrt2} = \sqrt2$: o lado companheiro é pequeno
demais. Se $x^2 < 2$, as desigualdades se invertem. De um jeito ou de
outro, $\sqrt2$ fica entre $x$ e $\frac2x$, então a média deles --- o
chute seguinte --- está mais perto do que o pior dos dois lados: o
retângulo de \omterm{def:g6:measure:area}{área} $2$ fica mais quadrado a cada passo.

\textbf{14.} $9 - 8 = 1$; \ $289 - 288 = 1$; \
$332\,929 - 332\,928 = 1$. Todo chute $\frac ab$ de Heron satisfaz
$a^2 - 2b^2 = 1$: a Parte I demonstrou que acertar o $0$ é impossível,
e essas frações erram o impossível por exatamente uma unidade --- o
mais perto que números inteiros conseguem chegar.

\textbf{15.} A reta numérica tem de conter números além das frações
--- comprimentos como $\sqrt2$, cujas escritas decimais correm para
sempre sem se repetir: os \emph{números reais}. A construção honesta
deles é assunto do volume do ensino médio e, com todo o rigor, dos
volumes de graduação.
\end{solution}
